\documentclass[11pt]{article}
\usepackage[margin=1in]{geometry}
\usepackage{setspace}
\usepackage{booktabs}
\usepackage{amsmath,amssymb}
\usepackage{graphicx}
\usepackage{svg}
\usepackage{float}
\usepackage[numbers,sort&compress]{natbib}
\usepackage[hidelinks]{hyperref}
\hypersetup{
  pdftitle={Household and Community Heterogeneity in Childhood Stunting in Ghana},
  pdfauthor={Gideon Ofosu Kyere},
  pdfsubject={Bayesian hierarchical analysis of childhood stunting in Ghana},
  pdfkeywords={childhood stunting, Bayesian hierarchical modelling, Ghana, household heterogeneity, community heterogeneity, WASH, maternal education}
}
\graphicspath{{../results/figures/}}
\svgpath{{../results/figures/}}

\title{Household and Community Heterogeneity in Childhood Stunting in Ghana}
\author{Gideon Ofosu Kyere, BSc\\\small Kwame Nkrumah University of Science and Technology, Kumasi-Ghana\\\small ORCID: 0009-0003-9848-8437\\\small Correspondence: 47 Mempeasem ST, GA-485-2566; kyereofosu2003@gmail.com}
\date{}
\begin{document}
\maketitle

\begin{abstract}
\noindent
Ghana's national stunting prevalence has fallen substantially, but the burden remains uneven. We conducted a cross-sectional secondary analysis of the 2022 Ghana Demographic and Health Survey to examine how much residual variation in childhood stunting lies between households and how much lies between survey communities. The analysis included 4,928 children aged 0--59 months in 3,545 households and 616 communities. Survey-weighted stunting prevalence was 17.4\% (95\% CI: 15.8--18.9). Bayesian hierarchical logistic models included household and community random intercepts, followed by WASH and maternal-education extensions. In the final model, the household SD was 1.23 compared with 0.39 at the community level; the corresponding household VPC and community ICC were 0.31 and 0.03. Boys had higher odds of stunting than girls (OR 1.50, 95\% posterior interval 1.24--1.82), as did children aged 24--35 months compared with those aged 0--5 months (OR 2.91, 2.06--4.21). Unimproved drinking-water source had OR 1.50 (1.15--1.99), while higher maternal education had OR 0.37 (0.20--0.65). Sensitivity analyses using alternative priors, survey weights, missing-data treatment, a spline for age, and more detailed WASH coding did not change the main pattern, although the water estimate was less precise in the weighted pseudo-posterior. Most unexplained clustering in these data was at the household rather than community level.
\end{abstract}

\noindent\textbf{Keywords:} childhood stunting; Bayesian hierarchical model; Ghana; Demographic and Health Survey; household heterogeneity; WASH; maternal education




\section{Introduction}

Ghana's national stunting prevalence has fallen, but childhood stunting remains unevenly distributed. In the 2022 Ghana Demographic and Health Survey (GDHS), about 17\% of children under five were stunted \citep{GSSICF2024}. The national estimate masks clear differences by region, household wealth, and place of residence. It also leaves an important question unanswered: after measured characteristics are taken into account, how much of the remaining variation is shared within households and how much is shared across survey communities?

The distinction matters because children in the same household share many immediate conditions, including food resources, caregivers, sanitation facilities, and economic circumstances. Households in the same survey cluster may also share markets, health services, infrastructure, and environmental exposures. If a model allows for community clustering but not household clustering, some of the similarity among children living together may be assigned to the wrong level.

Previous Ghanaian studies have documented socioeconomic, household, and geographical differences in stunting. Bayesian geostatistical analyses have described spatial variation \citep{AhetoDagne2021,AmoakoJohnson2022}, while multilevel studies have identified household and socioeconomic correlates \citep{AhetoEtAl2015,IddrisuGyabaah2023}. More recent work across 35 sub-Saharan African countries again points to differences by age, sex, maternal education, wealth, and contextual conditions \citep{MoloroEtAl2026}. A 2026 systematic review also found that WASH--stunting associations vary across exposure definitions and settings \citep{WijayantoEtAl2026}, and work using the 2022 GDHS has examined WASH and child nutrition among children born to young mothers \citep{OgumEtAl2026}. Against that background, our focus is narrower: we separate household from survey-community heterogeneity in the eligible under-five sample, then examine how that decomposition behaves as WASH and maternal education are added, using harmonized samples where needed.

We also examine whether the main findings depend on modelling choices that could affect interpretation, including prior scale, treatment of survey weights, missing maternal-education information, categorical versus spline-based age modelling, and binary versus more detailed WASH definitions.

\section{Methods}

\subsection{Data source and analytic sample}

This was a cross-sectional secondary analysis of the Household Member Recode from the 2022 GDHS, accessed under authorization from The DHS Program. Fieldwork was conducted from October 2022 to January 2023. The GDHS used a stratified two-stage cluster sample \citep{GSSICF2024}. Children were eligible for this analysis if they were aged 0--59 months, had slept in the sampled household on the night before the interview (\texttt{hv103 = yes}), and had a valid height-for-age z-score in \texttt{hc70}.

We defined stunting as height-for-age below $-2$ standard deviations from the WHO reference median, corresponding to \texttt{hc70 < -200} in the DHS recode. The resulting core sample contained 4,928 children in 3,545 households and 616 survey communities. For the hierarchical models, community was operationalized as the DHS survey cluster (\texttt{hv001}), and households were uniquely identified by their cluster and household identifiers (\texttt{hv001}, \texttt{hv002}).

\subsection{Covariates}

The core model included child sex, age, household wealth quintile, urban or rural residence, and administrative region. For the primary model, age was grouped as 0--5, 6--11, 12--23, 24--35, 36--47, and 48--59 months.

Drinking-water source (\texttt{hv201}) and sanitation facility (\texttt{hv205}) were grouped into improved and unimproved/no-facility categories using JMP-consistent facility-type definitions \citep{WHOUNICEFJMP2021}. These groupings refer to source or facility type; they do not represent the full JMP service ladders. Maternal education came from \texttt{hc61} and was classified as no education, primary, secondary, or higher. Maternal education was unavailable for 425 children, leaving 4,503 children for the complete-case maternal-education model.

\subsection{Descriptive analysis}

We used the DHS household-member weight (\texttt{hv005}) for descriptive prevalence estimates. Taylor linearization was used for standard errors and 95\% confidence intervals, with \texttt{hv021} as the primary sampling unit and \texttt{hv022} as the sampling stratum. The sum of the rescaled weights in our analytic sample was 4,292.97, which reproduces the official GDHS weighted height-for-age denominator of 4,293 after rounding.

\subsection{Bayesian hierarchical models}

For child $i$ in household $j$ and community $k$, we modelled
\[
Y_{ijk}\sim\text{Bernoulli}(p_{ijk}),
\]
with
\[
\text{logit}(p_{ijk})
=
\alpha+\mathbf{x}_{ijk}^{\top}\boldsymbol{\beta}+u_k+v_{jk},
\]
where $u_k\sim N(0,\tau_c^2)$ is the community random intercept and $v_{jk}\sim N(0,\tau_h^2)$ is the household random intercept. Both random effects were fitted with non-centred parameterizations.

We used $N(0,1.5^2)$ for the intercept, $N(0,1^2)$ for regression coefficients, and Half-Normal$(1)$ priors for the household and community standard deviations. Model 1 contained the core covariates together with household and community random intercepts. Model 2 added WASH. Model 3 added maternal education. We also refitted Model 2 on the 4,503-child maternal-education sample so that the Model 2--Model 3 comparison was not confounded by a change in sample composition.

For interpretation of the random effects, we used a latent logistic residual variance of $\pi^2/3$. The community ICC and household variance partition coefficient were
\[
ICC_c=
\frac{\tau_c^2}{\tau_c^2+\tau_h^2+\pi^2/3},
\qquad
VPC_h=
\frac{\tau_h^2}{\tau_c^2+\tau_h^2+\pi^2/3}.
\]
We also calculated median odds ratios (MORs) as
\[
MOR=\exp\left[\sqrt{2}\Phi^{-1}(0.75)\tau\right].
\]

\subsection{Computation, diagnostics, and sensitivity analyses}

Models were fitted with NUTS in PyMC. For the strengthened final Model 3 diagnostic run, we combined eight independent chains, each with 500 warmup and 500 retained draws, giving 4,000 retained posterior draws. We examined divergences, $\hat R$, effective sample size, BFMI, and tree depth, and used posterior predictive checks to compare replicated and observed prevalence.

We then challenged the main specification in several ways. We refitted Model 3 with tighter and wider priors; fitted a pseudo-posterior using survey weights rescaled to mean one; retained children with missing maternal education as a separate sensitivity category; replaced the age groups with a cubic B-spline; and replaced the binary WASH variables with more detailed water and sanitation categories. Predictive performance of harmonized Models 2 and 3 was compared with PSIS-LOO and Pareto-$k$ diagnostics. As an independent check on residual clustering and fixed-effect robustness, we also fitted logistic generalized estimating equation (GEE) models with exchangeable working correlation. The GEE working correlations were used only as sensitivity diagnostics and were not interpreted as Bayesian ICCs.

\subsection{Ethics and data access}

The 2022 GDHS protocol received ethical clearance from the Ethical Review Committee of the Ghana Health Service and from the ICF Institutional Review Board \citep{GSSICF2024}. The original survey used informed-consent procedures under its approved protocol. This secondary analysis used de-identified data obtained through authorized DHS access and involved no direct participant contact. We do not redistribute DHS microdata. The public reproducibility repository contains code and aggregate, non-identifying outputs only.

\subsection{Use of generative artificial intelligence}

OpenAI ChatGPT (GPT-5.6 Sol; accessed September 2026) assisted with language editing, manuscript organization, code drafting and review, and reproducibility documentation. It was not used to generate or alter the DHS microdata or the locked model outputs. The author reviewed and revised the resulting text and code, checked the reported numerical results against the saved outputs, and remains responsible for the analysis, interpretation, and integrity of the work.

\section{Results}

\subsection{Stunting prevalence and descriptive pattern}

Of the 4,928 children in the analytic sample, 927 were stunted. The survey-weighted prevalence was 17.4\% (95\% CI: 15.8--18.9). Boys had a higher weighted prevalence than girls (19.4\% vs 15.3\%), and prevalence was highest at ages 24--35 months (24.6\%). The wealth gradient was pronounced: 22.8\% of children in the poorest quintile were stunted compared with 7.5\% in the richest. Northern (29.6\%) and North East (29.3\%) had the highest regional estimates.


\begin{table}[H]
\centering
\caption{Selected characteristics of the analytic sample and survey-weighted stunting prevalence}
\label{tab:selected_descriptive}
\begin{tabular}{lrr}
\toprule
Characteristic & Unweighted $n$ & Weighted prevalence, \% (95\% CI) \\
\midrule
Overall & 4,928 & 17.4 (15.8--18.9) \\
Male & 2,529 & 19.4 (17.2--21.6) \\
Female & 2,399 & 15.3 (13.3--17.3) \\
Urban & 2,024 & 15.1 (12.6--17.5) \\
Rural & 2,904 & 19.6 (17.6--21.6) \\
Poorest wealth quintile & 1,601 & 22.8 (19.7--25.9) \\
Richest wealth quintile & 536 & 7.5 (4.7--10.3) \\
Improved water source & 3,904 & 16.4 (14.7--18.1) \\
Unimproved water source & 1,024 & 23.0 (19.0--27.0) \\
Improved sanitation facility & 2,219 & 13.8 (11.8--15.7) \\
Unimproved/no sanitation facility & 2,709 & 22.3 (19.9--24.6) \\
No maternal education & 1,468 & 24.4 (21.5--27.2) \\
Higher maternal education & 328 & 5.7 (2.4--8.9) \\
\bottomrule
\end{tabular}

\begin{minipage}{0.95\textwidth}
\footnotesize
\textit{Note.} Prevalence estimates use DHS sampling weights; confidence intervals use Taylor linearization with survey strata and primary sampling units. Maternal-education counts include only the observed education categories; 425 children had missing linked maternal education.
\end{minipage}
\end{table}

\begin{figure}[H]
\centering
\includesvg[width=0.90\textwidth]{stunting_by_region}
\caption{Survey-weighted prevalence of childhood stunting by region.}
\end{figure}

\subsection{Most residual variation was at the household level}

Adding a household random intercept changed the interpretation of the clustering structure. In Model 1, the community SD had a posterior median of 0.416 (95\% posterior interval 0.159--0.590), while the household SD was 1.066 (0.774--1.346). On the latent scale, the community ICC was 0.038 and the household VPC was 0.246. The corresponding MORs were 1.49 and 2.76.

The difference did not disappear in the fuller model. In the strengthened Model 3 run, the community SD was 0.389 and the household SD 1.234. The community ICC was 0.030, whereas the household VPC was 0.307. The MORs were 1.45 for communities and 3.24 for households. In other words, the broad household-versus-community contrast remained clearly apparent after the additional covariates were included.

\subsection{Fixed effects}

Several fixed effects were stable across the modelling sequence. In Model 3, boys had higher odds of stunting than girls (OR 1.50, 95\% posterior interval 1.24--1.82). Children aged 24--35 months had nearly three times the odds of those aged 0--5 months (OR 2.91, 2.06--4.21).

The socioeconomic gradient also remained. For children in the richest households, the posterior OR relative to the poorest group was 0.36 (0.20--0.64).

For WASH, the water result was clearer than the sanitation result. Unimproved drinking-water source had an OR of 1.50 (1.15--1.99). The corresponding sanitation estimate was 1.12 (0.86--1.46), leaving substantial posterior support for both modest harm and little or no independent association.

Maternal education showed a similar contrast across categories. Higher education was associated with lower odds of stunting relative to no education (OR 0.37, 0.20--0.65). The primary and secondary categories were much less distinct from the reference group: OR 1.02 (0.75--1.41) and 0.89 (0.68--1.14), respectively.


\begin{table}[H]
\centering
\caption{Selected posterior results from the strengthened final model}
\label{tab:main_bayesian_results}
\begin{tabular}{lcc}
\toprule
Parameter or comparison & Posterior median & 95\% posterior interval \\
\midrule
Male vs female OR & 1.50 & 1.24--1.82 \\
Age 24--35 vs 0--5 months OR & 2.91 & 2.06--4.21 \\
Richest vs poorest OR & 0.36 & 0.20--0.64 \\
Rural vs urban OR & 0.85 & 0.64--1.12 \\
Unimproved vs improved water OR & 1.50 & 1.15--1.99 \\
Unimproved/no vs improved sanitation OR & 1.12 & 0.86--1.46 \\
Maternal primary vs no education OR & 1.02 & 0.75--1.41 \\
Maternal secondary vs no education OR & 0.89 & 0.68--1.14 \\
Maternal higher vs no education OR & 0.37 & 0.20--0.65 \\
\addlinespace
Community SD & 0.39 & 0.07--0.61 \\
Household SD & 1.23 & 0.94--1.53 \\
Community ICC & 0.03 & 0.001--0.07 \\
Household VPC & 0.31 & 0.20--0.41 \\
Community MOR & 1.45 & 1.07--1.79 \\
Household MOR & 3.24 & 2.45--4.32 \\
\bottomrule
\end{tabular}

\begin{minipage}{0.95\textwidth}
\footnotesize
\textit{Note.} OR = odds ratio; SD = random-intercept standard deviation; ICC = intraclass correlation coefficient; VPC = variance partition coefficient; MOR = median odds ratio. Fixed-effect ORs and variance summaries come from the strengthened eight-chain Model 3 run.
\end{minipage}
\end{table}

\begin{figure}[H]
\centering
\includesvg[width=0.94\textwidth]{model3_posterior_or_forest}
\caption{Selected posterior odds ratios from the final maternal-education model.}
\end{figure}

\subsection{Diagnostics and sensitivity analyses}

The strengthened Model 3 run produced no divergences. BFMI ranged from 0.51 to 0.66 across chains, and no chain saturated the tree depth. Fixed-effect convergence was strong: the largest $\hat R$ among the reported fixed effects was about 1.003, and bulk effective sample sizes exceeded 2,200. Mixing was slower for the two variance parameters. The community and household SDs had $\hat R$ values of about 1.021 and 1.018, with bulk ESS of about 544 and 524. Both $\hat R$ values exceed the commonly used 1.01 guideline, so their exact interval endpoints warrant additional caution.

The posterior predictive distribution reproduced the observed overall prevalence and the prevalence in each of the six age groups. Changing the functional form of age made little difference to the other main coefficients: in the spline-age model, the OR was 1.49 (1.12--2.00) for unimproved water, 0.37 (0.20--0.67) for higher maternal education, and 0.35 (0.19--0.61) for richest versus poorest.

The more detailed WASH analysis helped locate the water signal. Surface-water use had an OR of 1.52 (1.10--2.08) relative to improved sources. The estimate for unprotected groundwater was 1.41 (0.90--2.21). For sanitation, neither open defecation [OR 1.23 (0.92--1.64)] nor other unimproved facilities [OR 0.93 (0.68--1.31)] was clearly separated from the improved-facility group.

The survey-weight pseudo-posterior was less precise for water: OR 1.38 (0.90--2.07). This did not reverse the direction of the estimate, but it weakened the strength of the evidence.

The independent GEE checks pointed in the same direction as the hierarchical model. The exchangeable working correlation was 0.154 when clustering by household and 0.024 when clustering by community. With household clustering, the GEE estimate for unimproved water was OR 1.40 (1.12--1.74); replacing age categories with a spline gave OR 1.39 (1.12--1.74). Under the detailed WASH coding, surface water had OR 1.43 (1.12--1.83), whereas unprotected groundwater remained less precise at OR 1.34 (0.93--1.94). These GEE quantities are sensitivity checks rather than substitutes for the Bayesian variance components.

Finally, observation-level PSIS-LOO did not distinguish the harmonized WASH and maternal-education models in a practically meaningful way. Model 3 improved ELPD by 0.92, with an SE of 3.38. No Model 3 observation had Pareto $k>0.70$; Model 2 had five, and neither model had $k>1$. Because individual children are left out while household and community effects are estimated from the remaining data, this comparison reflects conditional prediction within the observed hierarchy; it should not be read as performance for entirely new households or communities.

\section{Discussion}

The main result is the size of the household random effect. Once household and community clustering were modelled separately, residual variation between households was much larger than residual variation between survey communities. That ordering remained after WASH and maternal education were added.

This pattern is consistent with earlier Ghanaian work showing meaningful household-level variation in childhood malnutrition \citep{AhetoEtAl2015}, and it does not conflict with evidence of geographical heterogeneity from geostatistical studies \citep{AhetoDagne2021,AmoakoJohnson2022}. A DHS cluster captures a broader shared setting, whereas a household captures a more immediate one. Treating the cluster as the only higher level can blur these two sources of dependence.

The median odds ratios make the contrast easier to see: 3.24 at the household level compared with 1.45 at the community level in the final model. These values summarize the spread of the random effects; they are not causal effects of moving a child from one household to another. Even so, the size of the household term suggests that measured DHS covariates leave important shared household conditions unaccounted for. Food security, food allocation, maternal nutrition, caregiving, recurrent infection, and short-term economic shocks are plausible contributors.

The age pattern was clear: stunting was most common around the second and third years of life, and adjusted odds were highest at 24--35 months. Replacing the age groups with a spline left the other main coefficients almost unchanged, suggesting that the wealth, water, sex, and education estimates were not artifacts of the chosen age cut-points.

The wealth gradient was one of the most stable findings. Even after WASH and maternal education were in the model, children in the richest households had markedly lower odds of stunting than those in the poorest. Similar socioeconomic inequalities have been reported repeatedly in Ghana \citep{SeworJayalakshmi2024,OsborneEtAl2025}. The attenuation of the wealth coefficient after maternal education was added is unsurprising given the relationship between schooling and economic position, but these cross-sectional data cannot establish mediation.

The WASH results were less uniform. A binary unimproved-water indicator was associated with higher odds of stunting in the main model, but the weighted pseudo-posterior was more uncertain. The detailed coding added useful context: surface water showed the clearest signal, while unprotected groundwater was less precisely estimated. Sanitation, despite a large crude difference in prevalence, was weak after multivariable adjustment. Taken together, these results make a stronger case for saying that water source is associated with stunting in these data than for claiming an independent sanitation effect. They do not establish that changing water source alone would prevent stunting.

Higher maternal education was associated with substantially lower odds of stunting, whereas primary and secondary education were not clearly different from no education once the other covariates were included. Similar inverse associations between maternal schooling and child undernutrition have been reported elsewhere in Ghana and sub-Saharan Africa \citep{AhetoEtAl2015,TakeleEtAl2022,AgyenEtAl2024}. The pattern here should not be read as evidence of a precise educational threshold. Maternal education is also a marker for differences in health literacy, income opportunities, autonomy, healthcare use, and family circumstances that we cannot fully separate in this dataset.

The model comparison also deserves restraint. Adding maternal education was scientifically useful, but it did not materially improve observation-level predictive performance. The ELPD difference between harmonized Models 2 and 3 was small relative to its uncertainty. Because the comparison is conditional on the sampled hierarchy, it does not test prediction in entirely new households or communities. Maternal education is therefore retained for substantive reasons, not because Model 3 can be shown to predict better.

\subsection{Strengths and limitations}

The analysis has several practical strengths. The reconstructed descriptive sample agrees with the published GDHS after rounding for both the weighted denominator and national stunting prevalence. The models separate household from community clustering, use a harmonized sample when Models 2 and 3 are compared, and test the main findings under several alternative specifications rather than relying on a single fit.

The study also has important limitations. The GDHS is cross-sectional, so the odds ratios describe associations rather than causal effects. Several plausible determinants of child growth---including detailed diet, household food insecurity, maternal nutritional status, birth size, and repeated infection---were not available in the analysis. Some of that unmeasured information may be reflected in the household random effect.

Our WASH variables describe source and facility types, not the full JMP service ladders. Maternal education was missing for 8.6\% of the core sample, although a sensitivity analysis retaining those children did not materially change the main conclusions. The primary Bayesian likelihood does not fully incorporate the survey design; the weighted pseudo-posterior was therefore treated as a sensitivity analysis.

Finally, household variance is harder to estimate than the fixed effects. Many households contributed only one eligible child, and the two random-effect SDs mixed more slowly than the regression coefficients. The eight-chain run improved the effective sample sizes substantially and left the household-versus-community contrast unchanged, but exact variance-component intervals deserve more caution than the fixed-effect intervals.

\section{Conclusion}

The 2022 GDHS shows that childhood stunting in Ghana remains geographically and socioeconomically uneven, but the variation is not only between places. In these models, much more residual variation occurred between households within survey communities than between the communities themselves, and that pattern remained after WASH and maternal education were included.

The implication is that standard DHS covariates do not explain a substantial share of the clustering shared by children in the same household. Future work could examine that residual household component more directly using richer information on food security, diet, maternal health, infection, caregiving, and household shocks.


\section*{Data Availability Statement}
The 2022 Ghana DHS microdata are distributed by The DHS Program and cannot be redistributed by the author. Eligible researchers may request access directly from The DHS Program. The project repository contains code and aggregate, non-identifying outputs; it does not contain row-level DHS data.

\section*{Conflict of Interest Statement}
The author declares no conflict of interest.

\section*{Funding}
The author received no specific funding for this work.

\bibliographystyle{vancouver}
\bibliography{../thesis/references}
\end{document}
